\chapter{Modelos}

Este capitulo se centra principalmente en la investigación realizada para encontrar un modelo adecuado para la tarea del proyecto partiendo de un análisis inicial del problema hasta los posibles lenguajes candidatos. 

\section{Búsqueda de modelos}
\label{cap:modelos}
Tras realizar una búsqueda con las necesidades del proyecto se han elegido los siguientes candidatos:

\subsection{DeepSeek-Coder}
Se trata de un modelo entrenado con un total de 2 billones de tokens, en el que de esos tokens alrededor del 87\% es código y el resto son tokens en lenguaje natural en los idiomas inglés y chino. Ademas de que se nos ofrece el modelo con diferentes parámetros  de 1.3B (billones, entendidos como 1000 millones) a 33B, asi como un modelo base que sobretodo sirve para el autocompletado de código y otro \textit{instruct} (modelo capaz de entender instrucciones humanas) el cual esta entrenado para responder y resolver preguntas en lenguaje natural \cite{DeepSeek-coder}.

En el propio repositorio de Github nos ofrece una comparación con el resto de modelos open-source en los diferentes lenguajes de programación de una manera muy gráfica como vemos en la Figura \ref{fig:img-1}.

\imagen{img-1}{Grafico de comparación del rendimiento con respecto otros modelos open source basados en código \cite{DeepSeek-coder}}.

\subsubsection{Ventajas}
\begin{itemize}
	\item Es un modelo muy moderno, su primera versión es de noviembre de 2023.
	\item Como nos muestran en sus comparativas tanto el modelo con 7B y 33B parámetros tiene unos resultados muy buenos en todos los lenguajes superando a otros modelos como por ejemplo CodeLlama-34B en casi todos los lenguajes.
    \item Posee un modelo \textit{instruct} que esta adaptado tanto para recibir código como preguntas en lenguaje natural.
    \item Emplea licencia MIT
\end{itemize}
\subsubsection{Desventajas}
\begin{itemize}
	\item En los últimos meses se han realizado diferentes noticias sobre la seguridad del uso de DeepSeek con respecto a la protección de los datos.
    \item No dispone de tanta comunidad como otros modelos mas populares como CodeLlama o StarCoder
\end{itemize}

\subsection{CodeLlama}
Modelo perteneciente a meta derivado de Llama 2, ajustado para tareas de programación. Posee, ademas, el modelo base que permite autocompletado, un modelo especializado en python y otro que permite recibir instrucciones

\subsubsection{Ventajas}
\begin{itemize}
	\item Es un modelo que ha sido entrenado principalmente con conjuntos de datos de diferentes lenguajes de programación por lo que es ideal para este proyecto.
    \item Posee un modelo \textit{instruct} que esta adaptado tanto para recibir código como preguntas en lenguaje natural.
\end{itemize}
\subsubsection{Desventajas}
\begin{itemize}
    \item Dispone de una licencia mas restrictiva sin ser MIT.
    \item Es algo ''antiguo'', lo que puede resultar en que los resultados de inferencia no sean los mejores comparándolo con el resto de modelos.
\end{itemize}
\subsection{Llama 3}
Modelo perteneciente a meta permite comprender y generar texto mejor que sus predecesores con un buen razonamiento en código y conversación.
\subsubsection{Ventajas}
\begin{itemize}
	\item Es un modelo mas moderno, su fecha de lanzamiento fue el 18 de abril de 2024.
	\item Posee variantes del modelo desde 8B a 70B de parámetros 
    \item Posee un modelo \textit{instruct} que esta adaptado tanto para recibir código como preguntas en lenguaje natural.
    \item Al ser un lenguaje de meta y tener ya dos años en el mercado existen múltiples herramientas e implementaciones para el modelo que nos facilitan su implementación.
    \item Al no ser un modelo adaptado enteramente a código es capaz de entender mejor instrucciones en leguaje natural.
\end{itemize}
\subsubsection{Desventajas}
\begin{itemize}
    \item Dispone de una licencia mas restrictiva sin ser MIT.
    \item No es un modelo especializado en código lo cual si se piden tareas de código complejas podría generar respuestas menos optimas que otros lenguajes.
\end{itemize}
\subsection{StarCoder2}
StarCoder es un modelo de lenguaje especializado en código fuente, desarrollado por el proyecto BigCode, en  colaboración entre Hugging Face y ServiceNow Research. Es parte de la familia de modelos StarCoder / StarCoderBase, sucesores de SantaCoder, entrenados específicamente para tareas de desarrollo de software con múltiples lenguajes de programación y soporte para \textit{prompts} en lenguaje natural.
\subsubsection{Ventajas}
\begin{itemize}
	\item Es un modelo muy moderno, su primera versión es de noviembre de 2023.
	\item Como nos muestran en sus comparativas tanto el modelo con 15B y 3B parámetros tiene unos resultados buenos ademas de soportar multilenguaje.
    \item Los datos con los que ha sido entrenado pertenecen a repositorios de Github preseleccionados.
\subsubsection{Desventajas}
\end{itemize}
\begin{itemize}
    \item Dispone de una licencia mas restrictiva BigCode OpenRAIL-M v1.
\end{itemize}
\subsection{CodeGeeX2}
Modelo desarrollado por THUNLP (Tsinghua University NLP Group). Está diseñado para tareas de programación asistida por IA, con orientación multilingüe, tanto en lenguajes de programación como en lenguajes naturales.
\subsubsection{Ventajas}
\begin{itemize}
\item Tiene un tamaño de 6B de parámetros lo que permite despliegues locales o en servidores con GPU´s estandar.
\item Entrenado con lenguajes naturales incluido el español.
\item Dispone de licencia MIT.
\item Permite tareas de completado, explicación y depuración e código
\end{itemize}
\subsubsection{Desventajas}
\begin{itemize}
    \item Modelo menos conocido que CodeLlama o StarCoder. Si nos fijamos en las descargas del modelo de HuggingFace \cite{HuggingFace} tiene 153 con respecto a los otros dos modelos que tienen  4.997 y 172.744 respectivamente, lo que nos puede indicar que no es muy usado.
    \item No esta capacitado para autocompletado de código como otros modelos.
\end{itemize}